\begin{frame}{이번 주차 개요}
  ``당신은 가장 가까운 다섯 친구의 평균이다'' --- 이 한 줄이
  \kb{k-최근접이웃(kNN)}의 알고리즘 그 자체다. 거리를 직접 재는
  세 주인공: kNN(지도) · k-means(비지도) · GMM/EM(확률적 일반화).
  \vskip0.8em
  \begin{columns}[T]
    \begin{column}{0.55\textwidth}
      \textbf{이 주차를 마치면}
      \begin{itemize}\small
        \item kNN이 게으른(lazy)·비모수 방법임을, $k$를 키우며
          결정경계가 톱니바퀴 $\to$ 직선으로 변하는 과정을
          \kb{편향-분산 트레이드오프}로 설명
        \item 차원의 저주를 부피 공식 $p^{1/n}$과 이웃 간 거리 스케일로
          \textbf{정량적으로} 설명, 검증셋으로 $k$ 튜닝 코드 작성
        \item k-means 반복이 SSE를 단조감소시켜 수렴함을 수학으로 보임
        \item GMM이 왜 MLE로 바로 풀리지 않는지(우도가 \textbf{곱의
          합}), EM 한 라운드를 손계산
        \item 옌센 부등식 4줄로 ``우도가 절대 줄지 않는다'' 증명,
          $K$는 \kb{BIC}로 고르고 \kb{LDA}와 대비
      \end{itemize}
    \end{column}
    \begin{column}{0.43\textwidth}
      \textbf{네 차시 한 줄}
      \begin{itemize}\small
        \item \kb{4.1} 게으른 학습 · 거리함수·정규화 · 편향-분산
        \item \kb{4.2} 차원의 저주($p^{1/n}$, 동거리) · $k$ 튜닝 실습
        \item \kb{4.3} 수렴 증명 · 엘보우 · k-means++ · 구조적 한계
        \item \kb{4.4} GMM · EM · 옌센 부등식 · BIC · LDA
      \end{itemize}
      \vskip0.6em
      Week 3(생성 모델)이 회피한 \textbf{차원의 저주}를 이번 장은
      정면으로 맞는다. Week 5(SVM)로 가는 다리도 여기서 놓인다.
    \end{column}
  \end{columns}
\end{frame}
